\section{What remains open, and what was not checked}\label{sec:open}

\subsection{Open questions}

\begin{enumerate}[label=\arabic*.]
\item \emph{The global obligation.} Show that for every prime $p\equiv1,121,169,289,361,529\pmod{840}$ some shell is occupied. This is the conjecture itself, restated by Theorem~\ref{thm:shell}. The archive names it as the remaining obligation (p.~4). By Propositions~\ref{prop:jacobi} and~\ref{prop:groupbarrier}, any proof must produce, for some $R$, a shell whose subgroup $H_a$ contains $-1$ or $-4$. When $R$ is prime this means a prime factor of $(p+R)/4$ that is a non-residue modulo $R$. The proof must then go further, since that condition alone does not suffice; Proposition~\ref{prop:groupconverse} says when it does.
\item \emph{The census survivors.} $2.394\%$ of the final hard domain survives the carrier chain (§\ref{ssec:census}).
\begin{itemize}
\item The first half of the question asked before the referee pass is now answered. The census's carriers can never force the $2{,}759{,}345{,}613{,}832{,}500$ square rows, and no finite chain of such carriers can (Corollary~\ref{cor:censuslimit}).
\item Open: can the non-square survivor rows all be forced by carriers of the same kind, if more moduli are added?
\item Open: what kind of carrier, using the factorisation of $(p+R)/4$ and not only residues of $p$, could force the square rows?
\end{itemize}
\item \emph{The strict-record deficit-box conjecture} (§\ref{ssec:records}; archive §19.2). It holds at the two strict records with non-empty deficit below $8{,}803{,}369$ (reproduced here), and strict recordness cannot be dropped from it ($p=806{,}521$). A further test needs a strict record above $8{,}803{,}369$. There is none below $3\cdot10^9$ (Proposition~\ref{prop:records}).
\item \emph{Growth of the least class.} Is $h$ unbounded on the primes?
\begin{itemize}
\item Along the strict records up to $3\cdot10^9$, $h$ takes the values $1,2,3,6,8,15,27$.
\item Below $3\cdot10^9$ the next largest values are $21$, $20$ and $19$, at nine primes (Proposition~\ref{prop:records}).
\item Every record after $73$ lies in a hard class (Corollary~\ref{cor:hle2}).
\item The workbench states no conjecture on the growth of $h(p)$, and none is stated here.
\end{itemize}
\item \emph{The density of the shear failures.} What is the exact relative density, among the primes $p\equiv2\pmod3$, of the primes with a path of two shear edges? It is at least $1/24$, from the class $131\bmod180$, and at most $0.335$ (Remark~\ref{rem:sheardensity}). Below $10^7$ the proportion is $5.24\%$.
\item \emph{Type II universality} (Section~\ref{sec:notes}). Does every prime $p\equiv1\pmod{24}$ have a solution with $p$ dividing exactly two denominators? A yes proves the conjecture. Below $10^7$ every such prime has one in a shell $R\le107$, and only $p_*=8{,}803{,}369$ needs $107$.
\item \emph{The non-square classes of the notes' core} (Theorem~\ref{thm:termcore}). Is each of the $550$ non-square classes left open modulo $840\cdot11\cdot19\cdot23$ covered by finitely many single identities at some finite level? If one is not, then no identity sieve reduces the problem to the square core.
\item \emph{The quadratic-residue survivors} (Table~\ref{tab:survivors}). Are there infinitely many primes that satisfy every condition of §\ref{ssec:qrconditions}? A no would prove the conjecture outside a finite set of hard primes.
\item \emph{The $(3,4,\infty)$ covers at even levels.} At $D=2^k$ with $k\le5$ there are two orbits, of sizes $D^3/4$ and $3D^3/4$ (§\ref{ssec:es-s6}). Does this hold for every $k$? Is the orbit decomposition at a level $D=2^km$ with $m$ odd the product of those at $2^k$ and at $m$? At $D=12$ it is: the orbits have sizes $48$, $144$, $384$ and $1152$, the products of $16,48$ and $3,24$.
\end{enumerate}

\subsection{What this reader did not check}

\begin{itemize}
\item \emph{The archive.} Apart from the parts marked proved or reproduced in §\ref{ssec:archive}, its 619 pages were mapped, not audited. In particular, the side constructions (§§2, 4, 6--8, 9 after its opening, 21--23, the operator part of §24, the side material of §24.4, and §25) and the literature audits (§§11--14 and 17) were not compared with the papers they audit.
\item \emph{Lean.} The source archive contains 67 files with the extension \texttt{.lean} (seven of them \texttt{.axioms.lean}); they were not rebuilt. The archive's own account of what each Lean module checks is quoted where it matters (§\ref{ssec:census}).
\item \emph{Section~\ref{sec:notes}.} Several statements there rest on the reading passes and are marked so:
\begin{itemize}
\item the lattice and modular facts;
\item the mock theta and Virasoro computations;
\item the Busy-Beaver literature values;
\item the Cayley--Dickson identities;
\item the counts below $10^8$.
\end{itemize}
Salez's completeness statement was not proved here. The author's folder of drafts marked as overstated was not opened.
\item \emph{The census.} The enumerations of §\ref{ssec:census} were not re-run. Only the arithmetic consistency of the reported counts was checked, together with the base rows and the carrier sets used in Corollary~\ref{cor:censuslimit}.
\item \emph{The supplement.} The general proofs of the supplement's Sections 1--7 were not re-derived, except that of the no-two-shears theorem of Section 5 (Proposition~\ref{prop:twoshears}); their finite content was reproduced in part (§\ref{ssec:supplement}). Section 8, the auxiliary torus, was not checked.
\item \emph{The 124-page reader.} Only Proposition~\ref{prop:quartic} was checked.
\item \emph{The continuation.} Items 1, 5, 9--13, 15--20, 24, 25, 27 and 28 were located, not checked. So were item 4's enumeration, the deduction in item 23, and ES-04 to ES-08 beyond the ES-08 witness.
\item \emph{Novelty.} No literature search for novelty was made here. The workbench claims no novelty for the classical coordinates and identities it uses. The results marked as found by the referee pass (Propositions~\ref{prop:sheargraph}, \ref{prop:groupbarrier}, \ref{prop:groupconverse}, \ref{prop:fixeddivisor} and~\ref{prop:nosquare}, and Corollaries~\ref{cor:hle2} and~\ref{cor:censuslimit}) are short consequences of classical tools. Some may be known; none is claimed as new mathematics.
\end{itemize}
